% Double périodicité d'une onde sinusoïdale : longueur d'onde et période (L'essentiel, chapitre 17)
\begin{tikzpicture}
\begin{axis}[otMesures,otAxesBout,width=70mm,height=40mm,xmin=0,xmax=2.6,ymin=-1.4,ymax=1.6,
  axis x line=middle,xtick=\empty,ytick=\empty,xlabel={$x$},ylabel={$y$},name=g1,clip=false]
\addplot[otNavy,line width=1.1pt,domain=0:2.5,samples=120] {sin(deg(2*pi*x))};
\draw[otGreenDark,line width=0.7pt,{Stealth[length=1.5mm]}-{Stealth[length=1.5mm]}] (axis cs:0.25,1.25) -- (axis cs:1.25,1.25) node[pos=0.5,pcLegende,anchor=south,text=otGreenDark] {$\lambda$};
\draw[otMuted,line width=0.4pt,dash pattern=on 1.5pt off 1.2pt] (axis cs:0.25,1) -- (axis cs:0.25,1.3) (axis cs:1.25,1) -- (axis cs:1.25,1.3);
\node[pcLegendeG,anchor=north] at (axis cs:1.3,-1.45) {à une date donnée};
\end{axis}
\begin{axis}[otMesures,otAxesBout,width=70mm,height=40mm,xmin=0,xmax=2.6,ymin=-1.4,ymax=1.6,
  axis x line=middle,xtick=\empty,ytick=\empty,xlabel={$t$},ylabel={$y$},at={(g1.east)},anchor=west,xshift=12mm,clip=false]
\addplot[otOrange!85!black,line width=1.1pt,domain=0:2.5,samples=120] {sin(deg(2*pi*x))};
\draw[otGreenDark,line width=0.7pt,{Stealth[length=1.5mm]}-{Stealth[length=1.5mm]}] (axis cs:0.25,1.25) -- (axis cs:1.25,1.25) node[pos=0.5,pcLegende,anchor=south,text=otGreenDark] {$T$};
\draw[otMuted,line width=0.4pt,dash pattern=on 1.5pt off 1.2pt] (axis cs:0.25,1) -- (axis cs:0.25,1.3) (axis cs:1.25,1) -- (axis cs:1.25,1.3);
\node[pcLegendeG,anchor=north] at (axis cs:1.3,-1.45) {en un point donné};
\end{axis}
\end{tikzpicture}
\par\smallskip\legende{Périodicité spatiale (longueur d'onde $\lambda$) et périodicité temporelle (période $T$) : pendant une période, l'onde avance d'une longueur d'onde.}
